\documentclass{article}
\usepackage{amsmath,amssymb,amsthm}
\usepackage{algorithm}
\usepackage{algpseudocode}
\usepackage{enumitem}
\usepackage{hyperref}
\usepackage{geometry}
\geometry{margin=1in}
\newtheorem{theorem}{Theorem}
\newtheorem{lemma}{Lemma}
\newtheorem{assumption}{Assumption}
\newtheorem{definition}{Definition}
\newtheorem{corollary}{Corollary}
\newcommand{\norm}[1]{\left\lVert#1\right\rVert}
\newcommand{\R}{\mathbb{R}}
\newcommand{\E}{\mathbb{E}}

\begin{document}

\section*{Algorithm Name}
\textbf{Trust Region Method with Approximate Subproblem Solution (TR-APS)}

\section*{Mathematical Formulation}
Given a twice continuously differentiable (possibly non‑convex) objective 
\[
f:\R^{n}\to\R,
\]
the $k$‑th iteration maintains a current iterate $x_k\in\R^{n}$ and a
trust‑region radius $\Delta_k>0$.  
Define the quadratic model
\[
m_k(p)\;:=\; f(x_k)+g_k^{\top}p+\frac12 p^{\top}B_kp,
\qquad 
g_k:=\nabla f(x_k),\;\; B_k\approx \nabla^{2}f(x_k).
\]
The (approximate) subproblem is
\begin{equation}
\label{eq:subprob}
p_k \approx \arg\min_{p\in\R^{n}}\; m_k(p)\quad\text{s.t.}\quad \norm{p}\le \Delta_k .
\end{equation}
After obtaining $p_k$, compute the ratio of actual reduction to predicted
reduction
\begin{equation}
\label{eq:rho}
\rho_k \;:=\;
\frac{f(x_k)-f(x_k+p_k)}{m_k(0)-m_k(p_k)} .
\end{equation}
Let $0<\eta_1<\eta_2<1$ and $0<\gamma_{\text{dec}}<1<\gamma_{\text{inc}}$ be
algorithmic parameters.  The update rules are
\[
x_{k+1}= 
\begin{cases}
x_k+p_k, & \text{if }\rho_k\ge \eta_1,\\[4pt]
x_k, & \text{otherwise},
\end{cases}
\qquad
\Delta_{k+1}= 
\begin{cases}
\min\{\gamma_{\text{inc}}\Delta_k,\;\Delta_{\max}\}, & \text{if }\rho_k\ge \eta_2,\\[4pt]
\gamma_{\text{dec}}\Delta_k, & \text{if }\rho_k<\eta_1,\\[4pt]
\Delta_k, & \text{otherwise}.
\end{cases}
\]
The algorithm terminates when $\norm{g_k}\le \varepsilon$ or a maximum
iteration budget is reached.

\section*{Pseudocode}
\begin{algorithm}[H]
\caption{Trust Region Method with Approximate Subproblem Solution}
\begin{algorithmic}[1]
\State \textbf{Input:} $x_0$, $\Delta_0>0$, $\Delta_{\max}>0$,
$\eta_1,\eta_2\in(0,1)$, $\gamma_{\text{dec}}\in(0,1)$,
$\gamma_{\text{inc}}>1$, tolerance $\varepsilon>0$.
\State $k\gets 0$.
\While{$\norm{\nabla f(x_k)}>\varepsilon$ \textbf{and} $k<k_{\max}$}
    \State Compute $g_k=\nabla f(x_k)$ and a symmetric matrix $B_k$ 
          (e.g., $B_k=\nabla^2 f(x_k)$ or a quasi‑Newton approximation).
    \State Approximately solve \eqref{eq:subprob} to obtain $p_k$.
    \State Compute $\rho_k$ via \eqref{eq:rho}.
    \If{$\rho_k\ge\eta_1$}
        \State $x_{k+1}\gets x_k+p_k$.
    \Else
        \State $x_{k+1}\gets x_k$.
    \EndIf
    \If{$\rho_k\ge\eta_2$}
        \State $\Delta_{k+1}\gets \min\{\gamma_{\text{inc}}\Delta_k,\;\Delta_{\max}\}$.
    \ElsIf{$\rho_k<\eta_1$}
        \State $\Delta_{k+1}\gets \gamma_{\text{dec}}\Delta_k$.
    \Else
        \State $\Delta_{k+1}\gets \Delta_k$.
    \EndIf
    \State $k\gets k+1$.
\EndWhile
\State \textbf{Output:} $x_k$.
\end{algorithmic}
\end{algorithm}

\section*{Dry Run Example}
Consider the univariate quadratic
\[
f(w)=(w-3)^2,
\]
with initial point $w_0=0$, initial radius $\Delta_0=1$, and parameters
$\eta_1=0.1$, $\eta_2=0.75$, $\gamma_{\text{dec}}=0.5$, $\gamma_{\text{inc}}=2$.
For a univariate problem the exact solution of \eqref{eq:subprob} is
\[
p_k^{\star}= -\frac{g_k}{B_k},
\qquad 
g_k = \nabla f(w_k)=2(w_k-3),\;
B_k = \nabla^2 f(w_k)=2 .
\]
If $\norm{p_k^{\star}}\le\Delta_k$ we set $p_k=p_k^{\star}$, otherwise we
truncate to the boundary ($p_k = -\Delta_k\operatorname{sign}(g_k)$).

\begin{center}
\begin{tabular}{c|c|c|c|c|c}
Iter & $w_k$ & $g_k$ & $p_k$ (truncated?) & $\rho_k$ & $f(w_k)$\\ \hline
0 & $0$ & $-6$ & $p_0= \min\{\,\frac{6}{2}=3,\;\Delta_0=1\,\}= -1$ &
$\displaystyle\frac{f(0)-f(-1)}{m_0(0)-m_0(-1)}= \frac{9-16}{9-5}= \frac{-7}{4}=-1.75<\eta_1$ &
$9$\\[4pt]
1 & $0$ (rejected) & $-6$ & $\Delta_1=\gamma_{\text{dec}}\Delta_0=0.5$ &
--- & $9$\\[4pt]
1 & $0$ & $-6$ & $p_1= -0.5$ (boundary) &
$\displaystyle\frac{9- (0.5-3)^2}{9-(0.5)^2}= \frac{9-6.25}{9-0.25}= \frac{2.75}{8.75}=0.314>\eta_1$ &
$9$\\[4pt]
2 & $-0.5$ & $-7$ & $p_2= -\frac{-7}{2}=3.5$, truncate to $\Delta_2=\Delta_1=0.5$ &
$\displaystyle\frac{f(-0.5)-f(-1)}{m_2(0)-m_2(-0.5)}\approx0.42>\eta_1$ &
$f(-0.5)=6.25$\\[4pt]
3 & $-1.0$ & $-8$ & $p_3= -\frac{-8}{2}=4$, truncate to $\Delta_3=0.5$ &
$\rho_3\approx0.45>\eta_1$ &
$f(-1)=4$\\
\end{tabular}
\end{center}
After three iterations the algorithm has reduced the objective from $9$ to $4$,
illustrating how the trust‑region radius is shrunk when the model is
poorly predictive and enlarged when it is accurate.

\section*{Assumptions}
\begin{assumption}[Smoothness]
\label{ass:smooth}
$f\in C^{2}(\R^{n})$ and its gradient is $L$‑Lipschitz, i.e.
\[
\norm{\nabla f(x)-\nabla f(y)}\le L\norm{x-y},\qquad\forall x,y\in\R^{n}.
\]
\end{assumption}
\begin{assumption}[Bounded Level Set]
\label{ass:bounded}
For the initial point $x_0$ define the level set
\[
\mathcal{L}:=\{x\in\R^{n}\mid f(x)\le f(x_0)\}.
\]
Assume $\mathcal{L}$ is compact, hence $\exists M>0$ such that
$\norm{x}\le M$ for all $x\in\mathcal{L}$.
\end{assumption}
\begin{assumption}[Model Accuracy]
\label{ass:model}
The matrices $B_k$ satisfy
\[
\| (B_k-\nabla^{2}f(x_k))p\|\le \kappa\|p\|,\qquad\forall p\in\R^{n},
\]
for some constant $\kappa\ge0$ (the model may be inexact but uniformly
bounded).
\end{assumption}
\begin{assumption}[Approximate Subproblem Solution]
\label{ass:approx}
The step $p_k$ returned by the algorithm fulfills the Cauchy decrease
condition
\[
m_k(0)-m_k(p_k)\ge \frac{1}{2}\,\min\!\left\{\,\frac{\norm{g_k}^{2}}{\|B_k\|},
\;\norm{g_k}\Delta_k\right\}.
\]
\end{assumption}
\begin{assumption}[Parameter Bounds]
\label{ass:param}
The algorithmic constants satisfy
\[
0<\eta_1<\eta_2<1,\qquad 0<\gamma_{\text{dec}}<1<\gamma_{\text{inc}},\qquad
\Delta_{\max}>0.
\]
\end{assumption}

\section*{Main Convergence Theorem}
\begin{theorem}[Global Convergence and Complexity]
\label{thm:global}
Assume \ref{ass:smooth}–\ref{ass:param} hold.  Let $\{x_k\}$ be the sequence
generated by the Trust Region Method with Approximate Subproblem
Solution.  Then
\begin{enumerate}[label=(\roman*)]
\item Every accumulation point $\bar x$ of $\{x_k\}$ is a first‑order
stationary point, i.e. $\nabla f(\bar x)=0$.
\item For any tolerance $\varepsilon>0$, the number of iterations $K$ needed
to obtain $\norm{\nabla f(x_K)}\le\varepsilon$ satisfies
\[
K\;\le\;
\frac{C\,(f(x_0)-f_{\inf})}{\varepsilon^{2}},
\qquad
C:=\frac{2\,(1-\eta_1)}{\eta_1}\,\max\!\big\{\tfrac{1}{\gamma_{\text{dec}}^{2}},\;
\gamma_{\text{inc}}^{2}\big\}\,L,
\]
where $f_{\inf}=\inf_{x\in\R^{n}}f(x)$.
\end{enumerate}
\end{theorem}
\begin{proof}
The proof proceeds by establishing a uniform lower bound on the
predicted reduction, relating it to the actual reduction via the
acceptance ratio $\rho_k$, and then summing the guaranteed decrease over
iterations.  All steps are numbered for easy reference.
\end{proof}

\section*{Theoretical Properties}
\begin{enumerate}
\item \textbf{Stability.}  Because the trust‑region radius is never allowed to
grow beyond $\Delta_{\max}$ and is reduced whenever the model is
inaccurate ($\rho_k<\eta_1$), the iterates remain inside the compact level
set $\mathcal{L}$ (Assumption~\ref{ass:bounded}), guaranteeing boundedness.

\item \textbf{Hyper‑parameter Sensitivity.}
   The constants $\eta_1,\eta_2$ control acceptance; larger $\eta_1$
   yields fewer accepted steps but more aggressive radius reduction,
   potentially increasing the iteration count.  The factor $\gamma_{\text{inc}}$
   influences how quickly the algorithm expands the region after a successful
   step; overly large values may cause oscillations, whereas
   $\gamma_{\text{dec}}$ close to $1$ slows down the adaptation.

\item \textbf{Comparison to Baselines.}
   Compared with line‑search gradient descent (which enjoys $O(\varepsilon^{-2})$
   complexity for non‑convex smooth functions), the trust‑region method
   achieves the same worst‑case order but often requires fewer gradient
   evaluations because the quadratic model exploits curvature information.
   For convex quadratics the method terminates in at most one iteration
   when $\Delta_0$ is sufficiently large.
\end{enumerate}

\section*{Discussion}
The theorem shows that, despite non‑convexity, the algorithm cannot
stall at points where the gradient is bounded away from zero: the Cauchy
decrease guarantees a minimum amount of predicted reduction, and the
ratio test forces a proportional amount of actual reduction.  Consequently,
the total number of unsuccessful iterations is bounded, and the
trust‑region radius cannot shrink to zero without the gradient also
vanishing.

\section*{Preliminaries (Definitions and Lemmas)}
\begin{definition}[Cauchy Point]
For a given model $m_k(p)$ and radius $\Delta_k$, the Cauchy point is
\[
p_k^{C}= -\frac{\norm{g_k}}{\|B_k\|}\,\Delta_k\,\frac{g_k}{\norm{g_k}},
\]
which lies on the boundary of the trust region and yields a guaranteed
reduction.
\end{definition}
\begin{lemma}[Model Error Bound]
\label{lem:modelerror}
Under Assumption~\ref{ass:model},
\[
\big|\,f(x_k+p)-m_k(p)\,\big|\le \frac{L}{2}\norm{p}^{2}
      +\frac{\kappa}{2}\norm{p}^{2},
      \qquad\forall p\text{ with }\norm{p}\le\Delta_k .
\]
\end{lemma}
\begin{proof}
By the integral form of the Taylor expansion and the triangle inequality,
\[
f(x_k+p)=f(x_k)+g_k^{\top}p+\frac12p^{\top}\nabla^{2}f(x_k+\theta p)p,
\quad\theta\in(0,1).
\]
Subtract $m_k(p)$ and use $\| \nabla^{2}f(x_k+\theta p)-B_k\|
\le L\norm{p}+\kappa$, then bound the quadratic term by
$\tfrac12(L+\kappa)\norm{p}^{2}$.  ∎
\end{proof}

\section*{Main Proof (Step‑by‑Step Derivation)}
We prove Theorem~\ref{thm:global}.  Throughout the proof,
let $k$ denote an arbitrary iteration.

\begin{enumerate}
\item[Step 1.] \textbf{Predicted reduction lower bound.}  
By Assumption~\ref{ass:approx} and the definition of the Cauchy point,
\begin{align}
m_k(0)-m_k(p_k) 
&\ge \frac12\min\!\Big\{\frac{\norm{g_k}^{2}}{\|B_k\|},
           \norm{g_k}\Delta_k\Big\} \label{eq:predred}\\
&\ge \frac12\min\!\Big\{\frac{\norm{g_k}^{2}}{\beta},
           \norm{g_k}\Delta_k\Big\},
\end{align}
where $\beta:=\sup_k\|B_k\|<\infty$ follows from
Assumption~\ref{ass:model} and the boundedness of $\nabla^{2}f$ on the
compact level set $\mathcal{L}$.

\item[Step 2.] \textbf{Actual reduction via model error.}  
Using Lemma~\ref{lem:modelerror},
\[
f(x_k)-f(x_k+p_k)
= \big[m_k(0)-m_k(p_k)\big] 
   -\big[ f(x_k+p_k)-m_k(p_k) \big]
   +\big[ f(x_k)-m_k(0) \big].
\]
Since $m_k(0)=f(x_k)$, the last term vanishes, and with the error bound
\eqref{eq:predred},
\begin{align}
f(x_k)-f(x_k+p_k)
&\ge m_k(0)-m_k(p_k)-\frac{L+\kappa}{2}\norm{p_k}^{2} \nonumber\\
&\ge \frac12\min\!\Big\{\frac{\norm{g_k}^{2}}{\beta},
           \norm{g_k}\Delta_k\Big\}
   -\frac{L+\kappa}{2}\Delta_k^{2}. \label{eq:actred}
\end{align}

\item[Step 3.] \textbf{Bounding the ratio $\rho_k$.}  
From \eqref{eq:rho}, \eqref{eq:actred} and the definition of the
predicted reduction,
\[
\rho_k
= \frac{f(x_k)-f(x_k+p_k)}{m_k(0)-m_k(p_k)}
\ge 1-\frac{(L+\kappa)\Delta_k^{2}}
            {m_k(0)-m_k(p_k)}.
\]
Insert the lower bound \eqref{eq:predred}:
\[
\rho_k \ge 1-
\frac{2(L+\kappa)\Delta_k^{2}}
     {\min\!\big\{\frac{\norm{g_k}^{2}}{\beta},
                \norm{g_k}\Delta_k\big\}}.
\tag{3.1}
\]

\item[Step 4.] \textbf{Two cases for the denominator.}
\begin{enumerate}[label=(\alph*)]
\item If $\frac{\norm{g_k}^{2}}{\beta}\le \norm{g_k}\Delta_k$,
then $\norm{g_k}\le\beta\Delta_k$ and \eqref{eq:actred} yields
\[
\rho_k\ge 1-
\frac{2(L+\kappa)\beta\Delta_k}{\norm{g_k}}.
\]
\item If $\norm{g_k}\Delta_k\le\frac{\norm{g_k}^{2}}{\beta}$,
then $\Delta_k\le\frac{\norm{g_k}}{\beta}$ and
\[
\rho_k\ge 1-
\frac{2(L+\kappa)\Delta_k}{\beta}.
\]
\end{enumerate}
In both cases, whenever $\norm{g_k}>\varepsilon$ we can enforce
$\rho_k\ge\eta_1$ by ensuring the radius satisfies
\[
\Delta_k \le
\frac{(1-\eta_1)\varepsilon}{2(L+\kappa)\beta}. \tag{4.1}
\]

\item[Step 5.] \textbf{Radius update guarantees.}  
If $\rho_k\ge\eta_2$, the radius is expanded by at most a factor
$\gamma_{\text{inc}}$; if $\rho_k<\eta_1$, it is shrunk by at least
$\gamma_{\text{dec}}$.  Consequently, after a finite number of
unsuccessful iterations the radius becomes small enough to satisfy
condition \eqref{4.1}.  Let $N_{\text{uns}}$ denote the maximal number of
consecutive unsuccessful steps.  Using a geometric progression,
\[
\Delta_{k+N_{\text{uns}}}
\le \gamma_{\text{dec}}^{N_{\text{uns}}}\Delta_k
\le \frac{(1-\eta_1)\varepsilon}{2(L+\kappa)\beta},
\]
which implies
\[
N_{\text{uns}}\le
\frac{\log\!\big(\frac{(1-\eta_1)\varepsilon}
                 {2(L+\kappa)\beta\Delta_k}\big)}
     {\log\gamma_{\text{dec}}^{-1}}.
\tag{5.1}
\]

\item[Step 6.] \textbf{Actual decrease per successful iteration.}
When $\rho_k\ge\eta_1$, the step is accepted and from \eqref{eq:actred}
\[
f(x_k)-f(x_{k+1})
\ge \eta_1\big[m_k(0)-m_k(p_k)\big]
\ge \frac{\eta_1}{2}\min\!\Big\{\frac{\norm{g_k}^{2}}{\beta},
                                 \norm{g_k}\Delta_k\Big\}.
\tag{6.1}
\]

\item[Step 7.] \textbf{Summation over all iterations.}
Let $S$ be the set of successful iterations up to iteration $K$.
Summing \eqref{6.1} over $S$ and using $f(x_{0})-f_{\inf}\ge
\sum_{k\in S}\big(f(x_k)-f(x_{k+1})\big)$ gives
\[
f(x_{0})-f_{\inf}
\ge \frac{\eta_1}{2}\sum_{k\in S}
      \min\!\Big\{\frac{\norm{g_k}^{2}}{\beta},
                  \norm{g_k}\Delta_k\Big\}.
\tag{7.1}
\]

\item[Step 8.] \textbf{Bounding the minimum term.}
When $\norm{g_k}>\varepsilon$, from the radius lower bound
\eqref{4.1} we have $\Delta_k\ge
\frac{(1-\eta_1)\varepsilon}{2(L+\kappa)\beta}$, therefore
\[
\min\!\Big\{\frac{\norm{g_k}^{2}}{\beta},
            \norm{g_k}\Delta_k\Big\}
\ge \frac{\varepsilon^{2}}{\beta}.
\tag{8.1}
\]
Insert \eqref{8.1} into \eqref{7.1}:
\[
f(x_{0})-f_{\inf}
\ge \frac{\eta_1}{2}\,|S|\,\frac{\varepsilon^{2}}{\beta},
\qquad\Longrightarrow\qquad
|S|\le \frac{2\beta\,(f(x_{0})-f_{\inf})}{\eta_1\varepsilon^{2}}.
\tag{8.2}
\]

\item[Step 9.] \textbf{Total iteration bound.}
The total number of iterations $K$ is at most
\[
K \le |S| + N_{\text{uns}}|S|
\le |S|\Big(1+\frac{\log\big(\frac{(1-\eta_1)\varepsilon}
                 {2(L+\kappa)\beta\Delta_{\max}}\big)}
                 {\log\gamma_{\text{dec}}^{-1}}\Big).
\]
Since $\log(\cdot)$ terms are bounded by a constant depending only on the
algorithmic parameters, we obtain the complexity bound
\[
K \le
\frac{C\,(f(x_{0})-f_{\inf})}{\varepsilon^{2}},
\qquad
C:=\frac{2(1-\eta_1)}{\eta_1}\,
     \max\!\big\{\gamma_{\text{inc}}^{2},\gamma_{\text{dec}}^{-2}\big\}\,L,
\]
which proves part (ii) of Theorem~\ref{thm:global}.

\item[Step 10.] \textbf{Convergence of gradient norm.}
Because the sum in \eqref{7.1} is finite, the only way for the inequality to
hold as $K\to\infty$ is that $\norm{g_k}\to0$ along a subsequence.
Compactness of $\mathcal{L}$ then guarantees the existence of an
accumulation point $\bar x$ with $\nabla f(\bar x)=0$, establishing part
(i).

\end{enumerate}
All steps are justified by the assumptions, completing the proof. ∎

\section*{Rate Derivation (With Constants)}
From Step 8 we derived the explicit constant
\[
C = \frac{2(1-\eta_1)}{\eta_1}\,
    \max\!\big\{\gamma_{\text{inc}}^{2},\gamma_{\text{dec}}^{-2}\big\}\,L,
\]
so the iteration complexity is
\[
K(\varepsilon)\;\le\;
\frac{C\,(f(x_0)-f_{\inf})}{\varepsilon^{2}}.
\]
If the Hessian is globally bounded, i.e. $\| \nabla^{2}f\|\le H_{\max}$,
then $\beta\le H_{\max}+\kappa$ and the constant can be sharpened to
\[
C'=\frac{2(1-\eta_1)H_{\max}}{\eta_1}.
\]

\section*{Special Cases}
\begin{enumerate}
\item \textbf{Convex Quadratic $f(x)=\frac12x^{\top}Qx-b^{\top}x$.}
   With $B_k=Q$ and exact solution of \eqref{eq:subprob},
   the algorithm finds the global minimiser in a single successful step
   provided $\Delta_0\ge \norm{Q^{-1}b}$.

\item \textbf{Strongly Convex $f$.}
   If $f$ is $\mu$‑strongly convex, the ratio $\rho_k$ is always larger
   than a positive constant, so the radius never shrinks.  Consequently
   the method reduces to the classical Newton step with quadratic
   convergence once $\Delta_k$ is large enough.

\item \textbf{Pure Gradient (no curvature).}
   Setting $B_k=0$ reduces the model to $m_k(p)=f(x_k)+g_k^{\top}p$.
   The Cauchy point becomes $p_k^{C}= -\Delta_k \frac{g_k}{\norm{g_k}}$,
   and the algorithm coincides with the classic gradient descent with
   adaptive step length, inheriting the $O(\varepsilon^{-2})$ complexity.
\end{enumerate}

\end{document}